\documentclass{article}
\usepackage{amsmath,amssymb}

\title{Distributed Bump Hunting with Consensus-Coupled SGD}
\author{}
\date{}

\begin{document}
\maketitle

\section{Introduction}

The bump model suggests a distributed learning scheme that is close to
data-parallel SGD in implementation but very different in purpose.
Ordinary data parallelism tries to make many workers behave like a
single large worker: they process different batches, but they are
continually synchronized so that all replicas remain in essentially the
same part of parameter space. That is desirable when the goal is to
accelerate optimization of one trajectory.

Here we want something else. In a multi-bump landscape, different
regions of parameter space reveal different local structure. If every
worker is forced to stay near the same point, the system loses its main
distributed advantage: parallel discovery. We instead want many SGD
processes that remain genuinely separated in parameter space, so that
different workers can discover different bumps, ridges, or basins. At
the same time, we do not want those discoveries to remain private. Once
one worker finds a useful structure, the rest of the system should feel
that discovery and drift toward it.

This suggests a simple modification of data-parallel SGD: keep many
workers with their own parameters, let each worker take its usual local
stochastic gradient step, and then add a weak consensus step toward the
grand average of all workers. The average acts as a shared memory. A
worker that finds a new bump pulls the average toward that region; the
average then exerts a small force on the other workers, making the
discovery contagious without forcing immediate collapse to a single
point. The system is therefore neither independent multi-start SGD nor
standard synchronized SGD. It is a distributed search process with soft
communication through the ensemble mean.

The central claim of this paper is that such a system should inherit the
feature-discovery order of the bump model while exploring more
effectively than a single trajectory. Broad, low-rank bumps should be
found first by at least some workers. Once found, they bias the grand
average, which increases the probability that other workers also move
toward that region. Later, as workers accumulate around already-found
structure, the same mechanism should help the ensemble flow from broad
bumps toward sharper or intersecting bumps that would be difficult for a
single run to locate reliably.

\section{Consensus-Coupled SGD}

Suppose we have $m$ workers with parameters
$\theta_1, \dots, \theta_m \in \mathbb{R}^d$. Worker $i$ sees its own
mini-batch and computes a stochastic gradient $g_i(\theta_i)$ of a
common loss $L$. Let
\[
\bar{\theta} = \frac{1}{m} \sum_{j=1}^m \theta_j
\]
denote the grand average.\footnote{It might seem preferable to pull worker $i$ toward the leave-one-out average $\bar{\theta}_{-i} = \frac{1}{m-1}\sum_{j \neq i} \theta_j$ to avoid self-reinforcement. However, because $\theta_i - \bar{\theta} = \frac{m-1}{m}(\theta_i - \bar{\theta}_{-i})$, pulling toward the grand average with coupling $\gamma$ is mathematically identical to pulling toward the leave-one-out average with a slightly rescaled coupling $\gamma \frac{m-1}{m}$. For computational efficiency, we simply use the grand average.} The proposed update is
\begin{equation}\label{eq:update}
\theta_i^{t+1}
= \theta_i^t - \eta\, g_i(\theta_i^t)
- \gamma\,(\theta_i^t - \bar{\theta}^t),
\end{equation}
where $\eta > 0$ is the local learning rate and $\gamma > 0$ is a
consensus strength.

The first term is ordinary SGD. The second is the new ingredient: each
worker takes a small step toward the current center of mass of the
ensemble. If $\gamma$ is too large, the method degenerates into ordinary
data-parallel training because all workers collapse rapidly onto a
single path. If $\gamma$ is too small, the workers become nearly
independent and discoveries do not propagate. The interesting regime is
the intermediate one in which workers remain spatially separated for a
long time but still exchange information through the moving average.

Equation~\eqref{eq:update} can also be written as gradient descent on
the penalized objective
\[
\sum_{i=1}^m L_i(\theta_i)
+ \frac{\gamma}{2\eta} \sum_{i=1}^m \|\theta_i - \bar{\theta}\|^2,
\]
up to stochasticity and the dependence of $\bar{\theta}$ on all
particles. This makes the method look like a weakly coupled particle
system: local loss gradients pull each worker into nearby features,
while the quadratic coupling transmits discoveries across the ensemble.

\section{Why This Differs from Ordinary Data Parallelism}

Standard data-parallel SGD uses multiple workers to reduce gradient
variance and increase throughput, but it does not try to maintain
diversity. The workers are replicas. In effect, the system assumes that
all useful information can be aggregated at the gradient level around a
single current iterate.

The bump model suggests that this assumption is wrong when the loss
contains multiple separated structures. Far from a bump, the local
gradient may carry almost no information about that bump. A worker near
one bump and a worker near a different bump are not simply noisy
versions of the same state; they have access to different local
geometry. Forcing them to agree too quickly destroys that geometric
coverage.

Our proposal therefore treats spatial separation as a feature, not a
bug. The ensemble should occupy multiple regions at once. Discovery is
distributed, but sharing is centralized through $\bar{\theta}$. The
grand average is not meant to be the one true model at every moment. It
is a communication channel that accumulates weak evidence from all
workers.

\section{Qualitative Dynamics in a Bump Landscape}

The bump model gives a concrete picture for how the system should work.

\subsection*{Phase 1: independent exploration}

At initialization, the workers are scattered. Because the consensus term
is weak, each worker behaves roughly like an independent SGD trajectory
and can fall into a different broad basin. This is the phase where the
system gains coverage of parameter space.

\subsection*{Phase 2: discovery enters the average}

Suppose one worker encounters a broad low-rank bump first. That worker's
parameters begin to move coherently toward the bump center. As a result,
the grand average shifts slightly in the same direction. Every other
worker then feels a small drift toward the discovered structure, even if
its own local mini-batches have not yet revealed the bump clearly.

\subsection*{Phase 3: recruitment without immediate collapse}

If the coupling is weak enough, the other workers are not forced to jump
directly onto the first worker's path. Instead, they keep their own
local identities while being biased toward the promising region. Some
will join the same bump; others will continue to explore elsewhere.
This is the crucial regime in which discoveries spread through the
ensemble without destroying diversity.

\subsection*{Phase 4: flow toward intersections and sharper structure}

Once several workers are influenced by the same broad bump, their common
drift can carry parts of the ensemble toward intersections with other
bumps or toward sharper substructure inside the basin. In the language
of the original bump paper, the ensemble first finds broad low-rank
features and then collectively flows toward sharper or higher-rank
features. One worker may discover a side ridge, another may locate an
intersection, and the average can relay both discoveries back to the
rest of the system.

\section{The Averaging Trap and Soft Consensus}

To understand why this weak coupling is strictly necessary, consider the geometry of two workers discovering distinct features. We can simplify the problem by looking only at the 2D span connecting the centers of two separated bumps.

Suppose Worker A has successfully descended into the basin of Bump 1, and Worker B has descended into Bump 2. If the system uses standard synchronization (as in Federated Averaging or fully synchronous SGD), it forces both workers to jump directly to the grand average, $\bar{\theta} = \frac{1}{2}(\theta_A + \theta_B)$. Geometrically, this midpoint lies in the vast, empty plateau between the two bumps. By jumping to the average, the system instantly ejects both workers from their respective basins of attraction, destroying both discoveries simultaneously.

This ``averaging trap'' is why a soft consensus field is required. In our update rule~\eqref{eq:update}, workers do not jump to the average; they merely take a small gradient step towards it. This creates an elastic tension. The consensus term $-\gamma(\theta_i - \bar{\theta})$ pulls Worker A toward the empty plateau, but the local loss gradient $-\eta \nabla L(\theta_A)$ pulls it back down into the basin of Bump 1. 

As long as the coupling $\gamma$ is appropriately tuned, the local curvature of the bump will overpower the consensus pull. Worker A remains safely anchored in its basin, contributing its discovery to the shared memory $\bar{\theta}$, without being violently yanked into the void. This explains why consensus-coupled SGD can aggregate disjoint discoveries across a landscape where strict averaging would fail catastrophically.

\section{Historical Context: Swarms, MCMC, and the Curse of Dimensionality}

The idea of using a population of workers that explore independently and periodically share information is not new. It is the defining characteristic of several highly successful optimization and sampling families, most notably Particle Swarm Optimization (PSO) and Ensemble Markov Chain Monte Carlo (e.g., Parallel Tempering or Replica Exchange). In these algorithms, "particles" or "replicas" traverse the landscape, communicating their current fitness or state to the rest of the ensemble to prevent any single worker from becoming permanently trapped in a local minimum.

For decades, these methods have been the gold standard for complex, multimodal problems in physics and engineering. However, they are notoriously difficult to scale to the high-dimensional spaces of modern deep learning ($d \sim 10^6$ to $10^9$). The failure mode is geometric: both PSO and standard MCMC rely heavily on zero-order exploration (scalar fitness feedback) or random-walk proposals. In a million-dimensional space, the volume is so vast that randomly probing for a better state has an acceptance probability approaching zero. The "curse of dimensionality" guarantees that these swarms will wander aimlessly in the flat regions of the landscape, completely blind to narrow, high-rank structures.

Consensus-coupled SGD circumvents this curse by marrying the population-based communication of swarms with the directed, first-order efficiency of gradients. The local SGD step ($-\eta \nabla L$) is crucial: it provides an $n$-dimensional vector pointing directly into the basin of a bump. In the context of our model, the gradient step effectively collapses the million-dimensional search problem down to the intrinsic rank $r$ of the local bump. The worker easily slides down the $r$-dimensional ridge, solving the high-dimensional needle-in-a-haystack problem that defeats PSO.

Once the local gradient has anchored the worker in the low-dimensional structure, the consensus term ($-\gamma(\theta_i - \bar{\theta})$) acts as the swarm communication channel. It allows the ensemble to share the discovery of this $r$-dimensional subspace globally. Therefore, the success of EASGD and related methods on deep neural networks can be understood as a hybrid strategy: using gradients to defeat the curse of dimensionality locally, and using soft consensus to defeat local minima globally.

\section{Related Work: Elastic Averaging SGD}

The update rule proposed in~\eqref{eq:update} is mathematically equivalent to
Elastic Averaging SGD (EASGD), introduced by Zhang, Choromanska, and LeCun \cite{zhang2015deep}. In EASGD, local workers are coupled to a central center variable $\tilde{x}$ via a quadratic penalty $\frac{\rho}{2}\|x^i - \tilde{x}\|^2$, which acts as an elastic spring pulling workers toward the consensus mean.

Empirically, EASGD and its decentralized counterparts (e.g., D-PSGD \cite{lian2017can}) successfully reduce communication overhead and often achieve lower test error than fully synchronized methods like Downpour or standard All-Reduce SGD in non-convex landscapes. The original authors attributed this to the fact that the elastic force allows workers to fluctuate and explore different local optima rather than collapsing prematurely. However, in environments with fast interconnects or for tasks like speech recognition, highly synchronized methods or blockwise model-update filtering (BMUF) \cite{chen2016scalable} often remain competitive or superior in final accuracy.

The bump model provides a new geometric justification for \emph{why} EASGD works, and critically, suggests how it might be tuned to work better. Standard EASGD typically uses a fixed coupling constant $\alpha$. However, our bump analysis indicates that optimal discovery requires a phase transition: $\gamma$ must be extremely weak (or zero) early on to allow independent exploration of broad bumps, and then actively \emph{increased} (annealed upwards) to recruit the ensemble into sharper, high-rank structures discovered by individual pioneers. This schedule---weak early coupling followed by late consolidation---differs from the stationary elastic force typically employed.
\section{The X-Shaped Case}

The simplest test case is the X-shaped landscape formed by two rank-1
bumps crossing at roughly right angles. Suppose there are two workers.
One worker has drifted onto one ridge of the X, while the other has
drifted onto the orthogonal ridge. Their parameters are now different in
an important sense: the disagreement is not accidental mini-batch noise,
but evidence that the ensemble has discovered two distinct pieces of
structure.

This is exactly the regime in which naive consensus can fail. If the
coupling to the grand average is too strong, then each worker sees the
other worker's orthogonal displacement as a reason to move back toward
the center. The ensemble then collapses toward the intersection before
either worker has adequately explored its own ridge. In effect, the
system mistakes \emph{structured diversity} for \emph{error}.

What we want instead is asymmetric treatment of disagreement. Fast,
incoherent, worker-specific fluctuations should act like noise and
should not drag the whole system. But persistent directional
disagreement, especially disagreement that remains stable over multiple
steps, is precisely the information we want to preserve. In the X case,
the fact that two workers repeatedly occupy two nearly orthogonal
directions is a sign that the landscape contains two real ridges, not a
sign that the workers should immediately agree.

This means that coupling to the instantaneous mean
$\bar{\theta}^t = \frac{1}{m}\sum_j \theta_j^t$ is too crude. The raw
mean cannot distinguish between a transient deviation caused by
stochastic noise and a persistent deviation caused by discovery of a new
feature. To separate these, the shared state must be filtered.

The simplest modification is to replace the instantaneous center by a
slow center $c_t$:
\begin{align}
c_{t+1} &= (1-\beta)c_t + \beta \bar{\theta}^t, \\
\theta_i^{t+1} &= \theta_i^t - \eta\, g_i(\theta_i^t)
  - \gamma\,(\theta_i^t - c_t),
\end{align}
with a small tracking rate $\beta$. The point of $\beta$ is to make the
shared center respond only to disagreements that persist over time.
Short-lived fluctuations cancel out before they substantially move
$c_t$. But if one worker stays on one ridge and another stays on the
orthogonal ridge for many steps, then the center slowly moves toward the
intersection of those discoveries. The ridges are then communicated to
the rest of the population without forcing instant collapse.

In this picture, the orthogonal displacements of the two workers should
initially behave like \emph{private exploration}, not like a consensus
error to be erased. Only after that disagreement proves coherent over
time should it be promoted into the shared center. The center should
therefore represent \emph{durable discoveries}, not the average of every
momentary fluctuation.

The X-shaped example also clarifies the qualitative role of the two
coupling parameters:
\begin{itemize}
\item $\gamma$ controls how hard each worker is pulled toward shared
  structure.
\item $\beta$ controls how quickly the shared state believes that a new
  directional deviation is real.
\end{itemize}

For discovery, both should be modest. If $\gamma$ is too large, the
workers collapse before either ridge is explored. If $\beta$ is too
large, the center chases instantaneous deviations and converts noise
into spurious consensus. The useful regime is the one in which workers
can remain irrelevantly different in orthogonal directions for long
enough that the system can decide whether the disagreement is merely
noise or the signature of multiple bumps.

\section{Independent Draws and Consensus Speed}

The X-shaped discussion suggests a more quantitative design rule. The
consensus force should be strong enough to transmit genuine discoveries,
but weak enough that worker-specific motion in irrelevant directions is
treated like noise rather than immediate evidence for agreement. To make
that precise, we need a notion of when a worker has produced a new
\emph{effectively independent draw} in the irrelevant subspace.

Fix a worker near a ridge or bump and decompose its local coordinates
into a signal direction and an irrelevant orthogonal complement:
\[
\theta = (\theta_{\parallel}, \theta_{\perp}).
\]
The coordinate $\theta_{\parallel}$ is the part that carries the
discovery we want to preserve. The coordinate $\theta_{\perp}$ is the
part that should randomize quickly and behave like noise. Near a fixed
ridge, the stochastic dynamics of $\theta_{\perp}$ define a local Markov
process with transition operator $P_{\perp}$. Let
$1 = \lambda_1 > \lambda_2 \geq \lambda_3 \geq \cdots$ be its
eigenvalues. The second eigenvalue $\lambda_2$ controls the slowest
nontrivial mode of decorrelation.

A standard MCMC definition of the correlation time is
\[
\tau_{\mathrm{corr}} \asymp \frac{1}{1-\lambda_2},
\]
up to constants and logarithmic accuracy conventions. Equivalently, if
one wants an $\varepsilon$-accurate mixing time, one may write
\[
\tau_{\mathrm{mix}}(\varepsilon)
\asymp \frac{\log(1/\varepsilon)}{1-\lambda_2}.
\]
For the present paper, the exact convention is less important than the
timescale separation: the irrelevant coordinates become effectively
fresh after about $\tau_{\mathrm{ind}}$ steps, where
$\tau_{\mathrm{ind}}$ is any draw-spacing comparable to this mixing or
decorrelation time.

Now compare that with the consensus contraction. If the center were held
fixed for $\tau$ steps, the displacement of worker $i$ from that center
would shrink approximately like
\[
\theta_i^{t+\tau} - c \;\approx\; (1-\gamma)^{\tau}(\theta_i^t-c),
\]
ignoring the local SGD term for the moment. Thus the natural consensus
timescale is of order $1/\gamma$.

The design principle is then simple: do not let consensus erase worker
diversity faster than the irrelevant coordinates can produce a fresh
draw. If the worker-specific motion in $\theta_{\perp}$ has already
decorrelated after $\tau_{\mathrm{ind}}$ steps, then pulling on those
directions more slowly than that just averages noise. But if the signal
direction $\theta_{\parallel}$ remains coherent over that same window,
then averaging over $\tau_{\mathrm{ind}}$ steps keeps the signal while
washing out the irrelevant fluctuation.

This gives a concrete tuning rule. We would like the consensus to move a
worker only a moderate fraction of the way toward the shared center over
one independent-draw time. Your suggested target, ``move about halfway
toward the goal in the time it takes to get a new draw,'' corresponds to
\[
(1-\gamma)^{\tau_{\mathrm{ind}}} \approx \frac{1}{2}.
\]
For small $\gamma$, this yields
\[
\gamma \approx \frac{\log 2}{\tau_{\mathrm{ind}}}
\approx (\log 2)\,(1-\lambda_2).
\]
This is the right order of magnitude for the coupling strength. Larger
values force collapse before the irrelevant directions have randomized.
Much smaller values fail to propagate discoveries on the timescale at
which fresh exploratory information is being generated.

The same logic applies to the filtered center
\[
c_{t+1} = (1-\beta)c_t + \beta \bar{\theta}^t.
\]
If $1/\beta \ll \tau_{\mathrm{ind}}$, then the center tracks
instantaneous worker motion and mistakes noise for signal. If
$1/\beta \gg \tau_{\mathrm{ind}}$, the center becomes so sluggish that
real discoveries are communicated too late. Thus $\beta$ should also be
chosen on the scale of the independent-draw time, typically
\[
\beta \asymp \frac{1}{\tau_{\mathrm{ind}}},
\]
with a smaller constant than $\gamma$ when extra conservatism is
desired.

In the X-shaped case this is exactly the behavior we want. Orthogonal
motion along the irrelevant directions should decorrelate quickly, so it
looks like noise on the timescale $\tau_{\mathrm{ind}}$. But persistent
occupation of one ridge versus another should remain coherent across many
such windows. The filtered center then ignores the first kind of
disagreement and amplifies the second.

The broader lesson is that collapse should not be defined by raw worker
distance alone. It should be defined relative to a mixing time. Workers
have collapsed only when consensus contracts them together on a
timescale shorter than the time required for the irrelevant directions
to produce effectively independent draws. The spectral gap
$1-\lambda_2$ gives the right control knob for expressing that
comparison.

\subsection{Analytic Estimate of the Spectral Gap}

Because our data generating process~\eqref{eq:dgp} has a closed-form geometric definition, we can analytically estimate this spectral gap $1-\lambda_2$ and derive a concrete tuning rule for the coupling strength $\gamma$.

When a worker has successfully anchored onto the ridge of a bump with height $b$ and center $c=0$ (without loss of generality), it has found the optimal coordinate for the signal directions. We are interested in the dynamics of the orthogonal, irrelevant directions $\theta_{\perp}$. By definition, in these irrelevant directions, the shape matrix $A_i$ has zero eigenvalues. Therefore, the local loss function $f$ simplifies to the global decay term:
\[
f(\theta_{\perp}) \approx \frac{b}{1 + \|\theta_{\perp}\|^2}
\]
Since SGD is performing gradient \emph{ascent} to find the peak of this bump, the gradient in the irrelevant subspace is:
\[
\nabla_{\perp} f(\theta_{\perp}) = - \frac{2b \, \theta_{\perp}}{(1 + \|\theta_{\perp}\|^2)^2}
\]
Near the ridge ($\|\theta_{\perp}\| \approx 0$), we can linearize this gradient:
\[
\nabla_{\perp} f(\theta_{\perp}) \approx -2b \, \theta_{\perp}
\]
The SGD update step in these irrelevant directions is therefore a discrete-time Ornstein-Uhlenbeck (AR(1)) process:
\[
\theta_{\perp}^{t+1} \approx \theta_{\perp}^t + \eta \nabla_{\perp} f(\theta_{\perp}^t) + \xi_t
\]
\[
\theta_{\perp}^{t+1} \approx (1 - 2\eta b) \theta_{\perp}^t + \xi_t
\]
where $\xi_t$ is the mini-batch gradient noise. For an AR(1) process with decay coefficient $\rho = 1 - 2\eta b$, the eigenvalues of the transition operator $P_{\perp}$ are simply $\lambda_k = \rho^k$. Therefore, the second eigenvalue, which governs the slowest mode of decorrelation, is:
\[
\lambda_2 = 1 - 2\eta b
\]
The spectral gap is precisely $1-\lambda_2 = 2\eta b$. This gives us an exact, closed-form estimate for the independent-draw time:
\[
\tau_{\mathrm{ind}} \approx \frac{1}{1-\lambda_2} = \frac{1}{2\eta b}
\]
Plugging this spectral gap back into our design rule from the previous section ($\gamma \approx (\log 2)(1-\lambda_2)$), we obtain the optimal consensus coupling strength for the bump landscape:
\[
\gamma \approx (2 \log 2) \, \eta b \approx 1.38 \, \eta b
\]
This is a highly practical result. It demonstrates that the optimal consensus coupling $\gamma$ should be strictly proportional to the local learning rate $\eta$ and the height of the bump $b$. Importantly, because we are analyzing the orthogonal directions, the condition number or sharpest eigenvalues of the shape matrix $A_i$ do not enter this equation. The required coupling strength is dictated entirely by the global geometric ``gravity'' ($\|x\|^2$) holding the worker to the ridge, independent of how sharp the ridge itself is.

\section{Design Principles}

The bump picture suggests several operational principles.

\paragraph{Weak early coupling.}
The consensus strength $\gamma$ should initially be small relative to
the local SGD dynamics. Otherwise the workers collapse before they have
explored distinct regions. In the X-shaped case, this means the workers
must be allowed to remain separated along different ridges.

\paragraph{Broad initialization.}
The initial ensemble should cover many parts of space. If all workers
start nearly identical, the method reduces to ordinary data-parallel
SGD.

\paragraph{Late consolidation.}
After discovery, one may increase $\gamma$ or average more aggressively
so that the ensemble consolidates around the best found structures. The
early phase is for coverage; the later phase is for exploitation.

\paragraph{Persistent diversity.}
It may be useful to reserve some workers for continued exploration even
after the rest begin to consolidate. Otherwise the system risks
prematurely collapsing onto the first decent bump it finds.

\paragraph{Filtered consensus.}
Consensus should be applied to a temporally filtered shared state rather
than the raw instantaneous mean. This is the simplest way to let
durable discoveries influence the ensemble while treating fast
worker-specific fluctuations as noise.

\paragraph{Match consensus to mixing.}
The contraction rate $\gamma$ and center-tracking rate $\beta$ should be
chosen relative to the independent-draw time of the irrelevant
directions. A useful first rule is
$\gamma \approx (\log 2)/\tau_{\mathrm{ind}}$ and
$\beta \asymp 1/\tau_{\mathrm{ind}}$, where
$\tau_{\mathrm{ind}}$ is set by the decorrelation time or the inverse
spectral gap $1/(1-\lambda_2)$ of the local orthogonal dynamics.

\section{Predictions}

This view makes several qualitative predictions.

\begin{itemize}
\item On landscapes with multiple separated broad bumps, the method
  should discover useful basins faster than a single SGD trajectory.
\item Compared with standard synchronized data parallelism, the worker
  ensemble should exhibit larger pairwise distances early in training
  and better bump coverage.
\item The grand average should move before most workers individually
  reach a newly discovered bump, because it aggregates partial discovery
  from the first successful workers.
\item If $\gamma$ is increased too much, the system should lose its
  advantage and revert to single-trajectory behavior.
\item If $\gamma$ is too small, workers should remain diverse but fail
  to benefit from one another's discoveries.
\item In X-shaped or multi-ridge landscapes, a filtered center should
  preserve orthogonal worker discoveries longer than raw mean coupling
  and should reach the shared intersection more reliably.
\item The collapse threshold should be governed not only by worker
  distance but by the ratio between consensus time $1/\gamma$ and the
  independent-draw time $\tau_{\mathrm{ind}}$ of the irrelevant
  directions.
\end{itemize}

\section{Open Questions}

Several questions matter if this is to become a practical learning
system.

\begin{itemize}
\item What schedules for $\gamma$ best balance exploration and
  consolidation?
\item How should one estimate the local spectral gap or decorrelation
  time online from worker trajectories?
\item Should communication use the full mean $\bar{\theta}$, a robust
  average, or only a low-rank summary of the discovered directions?
\item Can one measure worker diversity and use it as a control signal
  for adaptive coupling?
\item What statistics best distinguish transient worker-specific noise
  from persistent directional disagreement that signals a real feature?
\item In realistic neural networks, which features correspond to the
  bumps that should be discovered in parallel, and at what layer or
  subspace should consensus operate?
\end{itemize}

The bump model does not answer these questions completely, but it gives
a sharp conceptual proposal: distributed learning should not merely
average many noisy copies of the same trajectory. It should maintain a
population of genuinely separated learners whose discoveries are shared
through a soft consensus field. That is the system suggested by the bump
landscape.

\bibliographystyle{plain}
\bibliography{distributed_bumps}

\end{document}
